\subsection{Teilaufgabe 2}
\subsubsection{Aufgabenstellung}
Es soll ein Programm erstellt werden welches zwei 2x2 Matrizen miteinander multiplizieren und addieren
kann. Wobei die Elemente der beiden Matrizen bei der Erzeugung zufällig initialisiert werden sollen. 

\subsubsection{Anforderungsdefinition}
\begin{enumerate}
	\item Addiere zwei 2x2 Matrizen.
	\item Multipliziere zwei 2x2 Matrizen.
	\item Initialisiere die Matrizen mit zufälligen zahlen.
\end{enumerate}

\subsubsection{Entwurf}
Es sollen zunächst 2 2x2 Matrizen erzeugt welche anschlie\ss end einer Methode übergeben werden,
welche die beiden Matrizen mit zufälligen Werten initialisiert. Danach sollen beide Matrizen einmal
ausgegeben werden. Nach der Ausgabe werden zwei Methoden aufgerufen, eine für die Multiplikation
und eine für die Addition, diese liefern jeweils eine neue Matrix zurück mit den ausgerechneten Werten.
Diese Matrix wird als MatrixC Temporär gespeichert und anschlie\ss end ausgegeben.
\begin{center}
	\includegraphics[width=0.45\textwidth]{uml/uml_c4_p2.pdf}
\end{center}

\subsubsection{Quellcode}
\paragraph{Matrizen.java}\
\lstinputlisting[language = Java , frame = trBL , escapeinside={(*@}{@*)}]{../chapter_04/src/main/java/chapter_04/Matrizen.java}

\subsubsection{Testdokumentation}
Das Programm hat nach dem Aufruf Zwei 2x2 Matrizen erstellt und initialisiert, anschlie\ss end 
miteinander Addiert und Multipliziert. Dabei wurde das Ergebnis mit einem Matrizen Rechner im Internet
verglichen und die Ergebnisse haben übereingestimmt.

\subsubsection{Benutzungshinweise}
Navigieren Sie in der Kommandozeile zum dem Ordner, wo sich die Java Datei befindet.
Danach führen sie "javac Matrizen.java\dq \space auf. Jetzt können Sie das Programm mit
"java Matrizen\dq \space starten. Nach dem das Programm gestartet ist, können Sie die
grö\ss e der Matrix angeben.

\subsubsection{Anwendungsbeispiel}
Nach dem Aufruf von java Matrizen, sollten wir nun folgendes sehen:
\begin{lstlisting}[frame = trBL , escapeinside={(*@}{@*)}]
[sebastian@laptop bin]$ java Matrizen
Matrix A:
70	50	
16	52	

Matrix B:
80	75	
11	33	

Addition von A und B:
150	125	
27	85	

Multiplikation von A und B:
6150	6900	
1852	2916
[sebastian@laptop bin]$  
\end{lstlisting}
